\documentclass[../../main.tex]{subfiles}

\graphicspath{{../../image/}{../../figure/}}

\begin{document}

% 原书物理页 7--8，印刷页码 VII--VIII。

\chapter*{第四版前言}
\phantomsection\label{chapter:front-preface}
\pdfbookmark[1]{第四版前言}{bm-front-preface}
\markboth{第四版前言}{第四版前言}

本书旨在对双曲守恒律理论作一广泛概述，并着重于它与经典连续介质物理之间的亲缘关系。本书的背景、范围与纲要在紧随本前言之后的\hyperref[chapter:front-introduction]{引言}中概述。本书初版于十五年前，第三版修订本于 2010 年问世。尽管如此，为了跟上该领域的新近进展，仍有必要准备这一经过大幅扩充与更新的新版。

面对研究在数量、多样性与技术复杂性上的爆炸式增长，本书原先那种百科全书式的雄心不得不有所节制。因此，若干重要的近期理论进展在此仅被略加触及，或只作勾勒。出于同样的原因，也无法呈现过去几年间如雨后春笋般涌现的大量各色应用。不过，经过更新的参考文献现已收录近两千条，仍为整个领域提供了一幅全景。

本书的根本目标在于促进守恒律双曲组的分析与连续介质物理之间的协同；这一目标在当下尤为切题，因为分析理论终于开始为迎接多空间维方程组所提出的挑战铺平道路。气体动力学的 Euler 方程组目前充当进入这一研究领域的入口。新版简述了该方向上的近期进展，并力求把弹性动力学方程组这一值得注意却未获应有重视的范例推到前台。

本版更加强调带耗散源、用以刻画松弛现象的双曲平衡律方程组理论。多空间维中关于经典解的那部分理论，在经大幅修订与扩充的\cref{chapter:ch05-熵与经典解稳定性}中阐述；而一维空间中的弱 $BV$ 解则在新增加的\cref{chapter:ch16-平衡律组bv解}中讨论。

原文中相当大的一部分已经重新组织，以使叙述更为精简、信息更为更新，并丰富例题的汇集。特别地，上一版的若干章通过增加新节而得以扩充，这些新节或对先前提出的问题作进一步阐发，或引入新的论题以供讨论。

\end{document}